% =====================================================================
% SEC 03 — Metodologia
% =====================================================================
\section{Metodologia}
\label{sec:metodo}

\subsection{Formulação da proposta pós-Recamán}

\subsubsection{A sequência de Recamán}
\label{sec:recaman-seq}

A sequência de Recamán $A005132$ \citep{Alekseyev2022RecamanCousins} é definida por
$a_0 = 0$ e, para $n\ge 1$,
\[
a_n \;=\; \begin{cases}
a_{n-1} - n, & \text{se } a_{n-1}-n > 0 \text{ e n\~ao visitado},\\
a_{n-1} + n, & \text{caso contr\'ario}.
\end{cases}
\]
Os primeiros termos são $0,1,3,6,2,7,13,20,12,21$.\footnote{Nota metodológica:
a sequência \textbf{não é injetora} para $n$ grande (revisita valores); qualquer
teste/benchmark que exija termo todos distintos é incorreto.}

\subsubsection{Offset de índice sobre o ranking}

Dado um ranking ordenado de $N$ candidatos e a sequência de Recamán
$a_0,a_1,\dots$, o índice candidato para o $m$-ésimo passo de diversificação é
\[
\label{eq:offset}
\mathrm{offset}_m \;=\; (1 + a_m) \bmod N .
\]
Para $N=8$, os offsets percorridos são $1,2,4,7,3,0,\dots$, cobrindo tanto o topo
quanto o fim do ranking com saltos curtos e longos alternados.

\begin{algorithm}[ht]
\caption{Diversificador pós-Recamán (determinístico)}
\label{alg:recaman}
\begin{algorithmic}[1]
\Require ranking $R$ de $N$ candidatos, orçamento $K$
\Ensure seleção $S$, com $|S|=\min(K,N)$, $|S|$ distinto, top-1 presente
\State $S \gets [\,]$; \Comment{seleção vazia}
\State $S.\mathrm{add}(R[0])$; \Comment{relevância primária nunca é perdida}
\State $\texttt{seen} \gets \{0\}$; $m \gets 0$;
\While{$|S| < \min(K,N)$}
  \State $p \gets (1 + a_m) \bmod N$;
  \If{$p \notin \texttt{seen}$}
    \State $S.\mathrm{add}(R[p])$; \Comment{``pula'' para posição $p$}
    \State $\texttt{seen} \gets \texttt{seen} \cup \{p\}$;
  \EndIf
  \State $m \gets m+1$;
\EndWhile
\Return $S$
\end{algorithmic}
\end{algorithm}

\subsubsection{Métrica de diversidade $\Div(\mathcal{S})$}

Sobre o conjunto selecionado $\mathcal{S}$, a diversidade média entre pares é
\[
\label{eq:diversity}
\Div(\mathcal{S})
\;=\; \frac{1}{|\mathcal{S}|\,(|\mathcal{S}|-1)}
\sum_{\substack{i\ne j\\ r_i,r_j\in\mathcal{S}}}
\bigl(1 - \Sim(r_i, r_j)\bigr),
\]
com $\Sim\in[0,1]$ uma similaridade normalizada. Valores próximos de $1$ indicam
alta diversidade; próximos de $0$, redundância.

\begin{quote}\small
\textbf{Insight de rigor.} Se $\Sim$ fosse medida sobre o \emph{texto bruto}
(coseno de embeddings), duas redações distintas da mesma fonte dariam $\Sim$
baixo, \emph{superestimando} a diversidade. Por isso a métrica vincula-se às
\emph{âncoras canônicas}: $\Sim$ é computada sobre identidade de
fonte/âmago (ex.: \texttt{source}). Isso garante \emph{diversidade efetiva}
(conteúdo realmente distinto), não \emph{variedade de redação}.
\end{quote}

\subsection{Formulação do HABD (diversificador híbrido anchor-blended)}

O achado do R458 — Recamán empata com o \emph{top-k} por diversificar
\emph{posição}, não \emph{âncora} — e a literatura de $\lambda$ adaptativo
\citep{Khan2026DFRAG,RezaeiDieng2025VendiRAG,GuoSanner2010PLMMR} motivam o
\textbf{HABD} (\emph{Hybrid Anchor-Blended Deterministic diversifier}), que
combina (a) o determinismo e o custo baixo do Recamán, (b) a dissimilaridade de
\emph{conteúdo} do MMR e (c) um $\lambda$ \emph{adaptativo por query}, sem LLM.
Formalmente, o HABD opera em três passos.

\subsubsection{Índice de mistura de âncoras ($\mu$)}

Dado o ranking de candidatos $R$ e a resolução de âncoras (Seção
\ref{sec:recaman-seq}), definimos o \emph{índice de mistura} sobre a janela de
topo de tamanho $W$:
\[
\label{eq:mu}
\mu(q) \;=\; \frac{
\mathclap{\bigl|\{ \text{\"ancoras distintas em } R[1..W] \}\bigr|}
}{
\min(W, |R|)
} ,
\qquad \mu \in [0,1].
\]
$\mu\approx 0$ indica ranking \emph{localmente monopolista} (poucas famílias no
topo); $\mu\approx 1$ indica ranking \emph{misto}. É uma medida barata e
determinística da \emph{estrutura} do ranking em relação às âncoras — a mesma
informação que, no DF-RAG \citep{Khan2026DFRAG}, é obtida por um julgador LLM.

\subsubsection{$\lambda$ adaptativo determinístico}

Para que a diversificação reaja ao monopolismo, mapeamos $\mu$ para $\lambda$ de
forma monotônica:
\[
\label{eq:lambda}
\lambda_{\mathrm H}(q) \;=\; \lambda_{\min} + \mu(q)
\left(\lambda_{\mathrm B} - \lambda_{\min}\right),
\]
com $\lambda_{\min}=0{,}3$ e $\lambda_{\mathrm B}=0{,}7$
($\mathrm H$ abrevia \emph{HABD} e $\mathrm B$, \emph{base}).
Quando $\mu\to 0$
(monopolista), $\lambda\to 0{,}3$ — a seleção passa a favorecer \emph{diversidade}
(penaliza fortemente âncoras repetidas). Quando $\mu\to 1$ (misto), $\lambda\to
0{,}7$ — privilegia-se \emph{relevância}, pois a diversidade já está
representada no ranking. Essa rotação entre relevância e diversidade é a ideia
central do \emph{over-retrieve + diversify} e da adaptação por tipo de consulta
presentes na prática moderna de RAG \citep{Khan2026DFRAG,RezaeiDieng2025VendiRAG}.

\subsubsection{Seleção MMR anchor-blended}

A escolha gulosa, a cada passo, maximiza:
\[
\label{eq:habdscore}
\mathrm{score}(d) \;=\; \lambda\,\mathrm{relev}(d) \;-\; (1-\lambda)\,
\underbrace{\max_{d_j\in S}\, \Sim_{\mathrm{anchor}}(d,d_j)}_{\mathclap{\text{0 se }d\text{ tem âncora nova}}} ,
\]
onde $\Sim_{\mathrm{anchor}}(d,d_j)=1$ se $d$ e $d_j$ têm a \emph{mesma âncora} e
$0$ caso contrário. A top-1 do ranking é sempre mantida. O efeito é duplo: (i)
uma \emph{âncora ainda não representada} no conjunto $S$ tem penalidade $0$ e,
portanto, prioridade — garantindo \emph{cobertura de fontes} determinística; (ii)
entre candidatos de \emph{mesma âncora}, vence o mais relevante. O custo é
$O(K\cdot |R|)$, idêntico ao do MMR, e o determinismo é preservado (sem seed;
desempates por ordem de ranking).

\subsection{Desenho do corpus-piloto}

Construímos um corpus determinístico com $A=6$ \emph{ângulos de perspectiva}
(famílias), cada um com $m=3$ documentos, totalizando $18$ documentos. Cada
ângulo associa-se a uma \emph{fonte} (\texttt{family\_$k$}) única e a um conjunto
de palavras-chave temáticas; o texto de cada documento recombina essas
palavras-chave de forma determinística. Dessa forma, a \emph{âncora canônica}
(coincidente com a \emph{fonte}) define a noção de ``família de perspectiva'' e
torna a métrica $\Div(\mathcal{S})$ coerente com a intuição de diversidade de
conteúdo.

Para cada query, um subconjunto de ângulos é designado como \emph{alvo} da
consulta; a query combina palavras-chave dos ângulos-alvo de forma \emph{equilibrada}
(ver os cuidados de validade adiante). O ranking de entrada é obtido por um
\emph{ranqueador} lexical-semântico determinístico (`ScientificRAG.retrieve`),
sobre o qual as estratégias de seleção operam.

\subsection{Protocolo de coorte}

Comparamos quatro estratégias de \emph{seleção pós-ranqueamento}, com o mesmo
orçamento $K$ por query:
\begin{enumerate}
  \item \textbf{\emph{top-k} (estado atual):} mantém os $K$ primeiros do ranking
        por relevância ($\mathrm{final\_score}$).
  \item \textbf{MMR:} seleção gulosa maximizando
        $\lambda\,\mathrm{relev} - (1-\lambda)\max \Sim(\cdot,\cdot)$, com
        $\lambda=0{,}7$ e similaridade por \emph{sobreposição de tokens}
        (Jaccard), determinístico e sem embeddings externos.
  \item \textbf{Recamán:} o \texttt{RecamanDiversifier} (Algoritmo~\ref{alg:recaman})
        sobre o ranking ordenado.
  \item \textbf{HABD:} o diversificador \emph{anchor-blended} determinístico com
        $\lambda$ adaptativo (Seção do HABD), com $\lambda_{\min}=0{,}3$ e
        $\lambda_{\mathrm B}=0{,}7$.
\end{enumerate}

As métricas avaliadas por seleção são: diversidade $\Div(\mathcal{S})$ (Eq.
\ref{eq:diversity} sobre âncoras), \emph{groundedness} (média dos
$\mathrm{final\_score}$ selecionados) e \emph{cobertura} (fração de ângulos-alvo
representados na seleção). O HABD registra, ainda, o $\lambda$ efetivo por query
(que varia conforme $\mu$), para verificação de que a rotação realmente ocorre.

\subsection{Cuidados de validade}

Para que o experimento seja \emph{válido} (não apenas executável), aplicamos
quatro correções metodológicas explícitas:

\begin{enumerate}
  \item \textbf{Query balanceada:} as palavras-chave da query são distribuídas
        proporcionalmente entre os ângulos-alvo, evitando que o 1.º ângulo domine
        por ordem de tokens.
  \item \textbf{Dedup por documento:} o ranqueador retorna múltiplos \emph{chunks}
        do mesmo documento; sem deduplicação, o topo é monopolizado por chunks do
        mesmo documento, tornando \emph{impossível} a qualquer diversificador
        espalhar. Deduplicamos por \texttt{doc\_id} antes da seleção.
  \item \textbf{Pool amplo:} coletamos um pool de candidatos (\texttt{poll\_size}
        $=40$ chunks) antes da deduplicação, para capturar mais de uma família.
  \item \textbf{Veredito anti-empate:} a hipótese H2 exige ganho \emph{estrito} de
        diversidade; empate ($\Div_{\text{Rec}}\le \Div_{\text{top-k}}$) é tratado
        como \emph{não-demonstração}, não como sucesso.
\end{enumerate}

\subsection{Hipóteses}

Formalizamos as hipóteses a testar, usando subscritos $t$ (top-$k$),
$r$ (Recamán), $m$ (MMR) e $h$ (HABD) para abreviar $g$ e $\Div$. Para o esquema
posicional de Recamán:
\[
\begin{aligned}
\textbf{H2:}\quad & \Div(S_r) > \Div(S_t)
  \;\wedge\; \tfrac{g_t - g_r}{g_t} \le 0{,}05,
\end{aligned}
\]
onde $g$ denota \emph{groundedness} médio. A segunda condição garante que o
ganho de diversidade, se houver, não degrada relevância além de uma tolerância de
$5\%$.

Para o diversificador híbrido HABD:
\[
\begin{aligned}
\textbf{H3:}\quad & \bigl[\Div(S_h) > 0{,}5\bigr]
  \;\wedge\; \bigl[\Div(S_h) \ge \Div(S_m)\bigr]\\[-1pt]
&\qquad \wedge\; \tfrac{g_t - g_h}{g_t} \le 0{,}05,
\end{aligned}
\]
em que a primeira condição exige superar o \emph{empate} top-k/Recamán, a segunda
exige não ficar atrás do MMR em diversidade, e a terceira impõe a mesma tolerância
de relevância. Ambas as hipóteses são \emph{refutáveis}: se a evidência não as
sustentar, registramos \texttt{refuta\_H2} / \texttt{refuta\_H3} sem maquiar. O
HABD \emph{pode} superar em diversidade e cobertura (factual, reportável) e ainda
assim \emph{refutar} H3 por não respeitar a tolerância de relevância — as duas
dimensões (ganho observável vs. critério rígido) são explicitamente separadas.
